% =========================================================
% DISCUSIÓN — versión RUCP (~1,200 palabras)
% =========================================================

\section{Discusión}
\label{sec:discusion}

\subsection{Interpretación de los Hallazgos}

El hallazgo central de este estudio ---que los circuitos rurales de Uruguay exhibieron tasas de defección de CA al FA 21{,}8 pp superiores a los urbanos--- contradice la narrativa convencional de conservadurismo rural \citep{gonzalez1991electoral, rodden2019s}. Tres patrones interrelacionados reconfiguran la comprensión del balotaje de 2024: (1) la brecha urbano-rural es robusta y asimétrica (afecta a CA y PC pero no al PN); (2) la heterogeneidad departamental supera en magnitud a cualquier estratificación binaria (rango de 89.6 pp entre Tacuarembó y Río Negro); (3) la nula estabilidad geográfica entre 2019 y 2024 ($r < 0.13$) indica reconfiguración estructural, no continuidad.

El contraste entre la estabilidad del PN y la volatilidad de CA y PC en el estrato rural sugiere que el mecanismo explicativo está en la solidez de los vínculos partidarios \citep{buquet2015uruguayan}. Los votantes del PN en zonas rurales mantienen una identidad organizacional con el partido que históricamente se enraíza en el interior agropecuario \citep{gonzalez1993structuring}; los votantes de CA, en cambio, fueron atraídos en 2019 por un liderazgo personal y posicionamiento anti-establishment que no genera el mismo nivel de arraigo institucional \citep{perezbentancur2019frente}. Ante la elección binaria del balotaje \citep{birch2003tworound} y el deterioro económico rural (sequía, precios de commodities, costos de producción), estos votantes de bajo arraigo defeccionaron masivamente.

\subsection{Implicaciones Teóricas}

Los resultados tienen tres implicancias para la literatura sobre comportamiento electoral comparado. \textit{Primera}: la dicotomía rural-conservador / urbano-progresista, documentada en democracias industrializadas \citep{rodden2019s}, no se transfiere automáticamente al contexto de segundas vueltas en presidencialismos latinoamericanos \citep{chasquetti2004democracia}. El voto en segunda vuelta puede estar más determinado por la evaluación retrospectiva del desempeño gubernamental \citep{downs1957economic} que por la identidad geográfica, especialmente cuando el partido del votante no está en el balotaje. \textit{Segunda}: la ``paradoja de la coalición cohesiva derrotada'' ilustra que en sistemas con elección de dos rondas, la capacidad de retener votos en la segunda vuelta no basta para ganar si el caudal electoral de la primera vuelta se erosionó entre elecciones. La victoria requiere crecer o mantener la base en la primera vuelta, no solo optimizar la transferencia en la segunda \citep{birch2003tworound}. \textit{Tercera}: la nula estabilidad geográfica de las defecciones entre 2019 y 2024 sugiere que la heterogeneidad electoral en Uruguay refleja condiciones coyunturales locales más que identidades geográficas duraderas ---un hallazgo relevante para la literatura sobre sistemas de partidos institucionalizados \citep{mainwaring1995party, luna2014segmented}.

\subsection{Limitaciones}

Cinco limitaciones merecen reconocimiento explícito. \textit{Primera}: el modelo de complete pooling estima una única $\mathbf{T}$ para todos los circuitos, ignorando la heterogeneidad dentro de los estratos. El modelo jerárquico con partial pooling por departamento (implementado en \texttt{HierarchicalEI}) es la extensión natural y constituye trabajo futuro prioritario. \textit{Segunda}: el problema de inferencia ecológica es inherentemente no identificado: múltiples matrices de transición individuales son consistentes con los mismos datos agregados \citep{king1997, greiner2009rxc}. Las estimaciones son consistentes con los datos pero no pueden verificarse sin encuestas de seguimiento. \textit{Tercera}: la likelihood Normal es una aproximación ad-hoc a la Multinomial; el modelo con verosimilitud DirichletMultinomial \citep{mosimann1962compound} confirma la dirección de los hallazgos y es preferible estadísticamente, pero requiere más tiempo computacional. \textit{Cuarta}: la clasificación urbano/rural proviene de datos administrativos que pueden no capturar la continua gradación entre entornos periurbanos y rurales profundos. \textit{Quinta}: no se evalúan los supuestos de invarianza de la matriz de transición dentro de cada circuito, ni la independencia entre circuitos.

\subsection{Direcciones de Investigación Futura}

Cuatro extensiones son prioritarias. (1) \textbf{Modelo jerárquico con covariables}: incorporar turnout, tamaño del circuito, proporción de CA en primera vuelta y clasificación urbana como predictores de la concentración del prior DirichletMultinomial, para reducir la anchura de los ICs y hacer interpretables los efectos de las covariables. (2) \textbf{Microdatos del Censo 2023}: cuando los microdatos del INE estén disponibles, integrar covariables socioeconómicas (educación, edad mediana, índice de ruralidad) para modelar heterogeneidad a nivel departamental. (3) \textbf{Comparación con eiCircles}: triangulación metodológica con el paquete \texttt{eiCircles} \citep{brown1986aggregate} vía la función \texttt{BPF()}, que implementa el modelo de Forcina et al.\ (2012). (4) \textbf{Encuestas de boca de urna}: la validación con datos individuales permitiría verificar las estimaciones ecológicas y testear mecanismos causales para la rebelión rural.
